% ECONOMICS_PROBLEM: {"id": "E43-556", "title": "Measuring the neutral interest rate", "jel_code": "E43", "source_id": "OP-0245"}
% !TeX program = xelatex
\documentclass[11pt,a4paper]{article}
\usepackage[margin=23mm]{geometry}
\usepackage{fontspec}
\setmainfont{Latin Modern Roman}
\usepackage{amsmath,amssymb,mathtools}
\usepackage{xcolor,enumitem,booktabs,longtable,array,xurl,hyperref}
\definecolor{muted}{HTML}{53636C}
\hypersetup{colorlinks=true,urlcolor=blue,linkcolor=blue}
\setlength{\parindent}{0pt}
\setlength{\parskip}{5pt}
\newcommand{\R}{\mathbb R}
\newcommand{\E}{\mathbb E}
\newcommand{\N}{\mathbb N}
\newcommand{\ind}{\mathbf 1}
\newcommand{\Var}{\operatorname{Var}}
\newcommand{\Cov}{\operatorname{Cov}}
\newcommand{\supp}{\operatorname{supp}}
\DeclareMathOperator*{\argmin}{arg\,min}
\DeclareMathOperator*{\argmax}{arg\,max}
\newcommand{\entrymeta}[3]{{\sffamily\small\color{muted}\textbf{\texttt{#1}}\quad|\quad #2\par #3\par}}
\newcommand{\Problemref}[1]{\texttt{#1}}
\newcommand{\JELref}[2]{\texttt{#2} (JEL \texttt{#1})}
\begin{document}
\section*{Measuring the neutral interest rate}
\textbf{Problem ID:} \texttt{E43-556}\quad \textbf{JEL:} \texttt{E43}\par
% BEGIN ECONOMICS STATEMENT
\label{problem:OP-0245}
{\sffamily\small\color{muted}Original subject group: Inflation and monetary policy\par}
\label{definitions:problem:30007526}
\textbf{Audit classification:} Background; proposed extension. Metadata and abstract are verified; the exact state-space identification target has no verified published open statement. Verification concerns the cited version, not first priority or proof that this exact editorial specification remains unsolved on October 8, 2026.
\subsection*{Definitions and precise research target}
\paragraph{Editorial mathematical specification.} In a specified small macroeconomic state-space model, let \(x_t\) be the output gap, \(r_t\) the ex ante real policy rate, and \(r_t^*\) the latent contemporaneous neutral rate. Assume
\[
x_t=E_tx_{t+1}-\sigma^{-1}(r_t-r_t^*)+d_t,\qquad
\pi_t=\beta E_t\pi_{t+1}+\kappa x_t+u_t,
\]
where \(\sigma,\kappa>0\), \(\beta\in(0,1)\), \(d_t\) is an additional demand disturbance, and \(u_t\) a cost disturbance. Specify observed proxies, measurement errors, transition laws, and survey information. Define the long-run natural rate \(\bar r\) as the steady-state flexible-price real rate under stated demographic, productivity, and fiscal fundamentals; it need not equal \(r_t^*\).
Determine whether \(r_t^*\), \(\bar r\), and their difference are identified from the proposed information set. In particular, exhibit or eliminate observational equivalence between shifts in \(r_t^*\), demand disturbances, and revisions to potential output. Develop confidence sets accounting for weak identification and model uncertainty, then assess their real-time stability across data vintages and policy regimes. A smooth filtered latent series is insufficient evidence of identification. The target concerns reliable distinctions between current inflation-neutral conditions and long-run equilibrium in this declared model family, rather than measurement of an observable universal interest-rate constant.

The displayed IS equation identifies at most \(r_t^*+\sigma d_t\) without further restrictions: replacing \(r_t^*\) by \(r_t^*+a_t\) and \(d_t\) by \(d_t-a_t/\sigma\) leaves that equation unchanged. Thus a transition restriction, excluded observable, or externally identified demand shock is indispensable. Specify the state-space observation equation \(z_t=H_\theta s_t+v_t\), transition \(s_{t+1}=A_\theta s_t+\eta_{t+1}\), initial state law, and covariances of \((v,\eta)\), with \(r_t^*\) a declared coordinate. In this specification \(\bar r\) is the flexible-price steady-state rate; it is not automatically the time average of a filtered \(r_t^*\). Identification means that all parameter/state laws generating the same observable law give the same target; a model-imposed smoothing assumption may identify a series conditionally without validating the assumption.
\subsection*{Variants and assumption changes}
The following are \emph{editorial variants}, not additional published open conjectures.
\begin{enumerate}
\item \emph{Restricted demand shock.} Set \(d_t=0\) or impose a known covariance/transition law for it. Compare the identified rate set with the unrestricted-demand case and report how much identification comes from the restriction.
\item \emph{Open-economy natural rate.} Add international asset pricing and global saving determinants. Compare a country-specific flexible-price rate with the world rate under imperfect financial integration, keeping neutral-rate concepts distinct.
\end{enumerate}
\subsection*{References and source audit}
\begin{itemize}
\item Maurice Obstfeld. \emph{Natural and Neutral Real Interest Rates: Past and Future}. 2023. Locator: Primary NBER indexed metadata and abstract; original December 2023, revised March 2025; direct fulltext unavailable. Primary indexed metadata/abstract checked; direct page and PDF access failed. \url{https://www.nber.org/papers/w31949}.
\end{itemize}
Metadata and abstract are verified; the exact state-space identification target has no verified published open statement.
\begin{enumerate}\raggedright
\item\hypertarget{definitions:source:30007526:1}{} Maurice Obstfeld. 2023. \emph{Natural and Neutral Real Interest Rates: Past and Future}. Primary NBER indexed metadata and abstract; original December2023, revised March2025; direct fulltext unavailable.
\url{https://www.nber.org/papers/w31949}.
DOI: \url{https://doi.org/10.3386/w31949}.
\end{enumerate}

\subsection*{Common conventions: Source mathematical conventions}

These conventions apply to each entry unless explicitly replaced.

\textbf{Models and identification.} A primitive model \(m\) specifies all probability laws, feasible choices, information, equilibrium and measurement rules used by its target. The admissible class \(\mathcal M\) is fixed independently of outcomes. If \(P_O(m)\) is its observable probability law and \(\tau(m)\) the target, its identified set at observable law \(P\) is
\[
 \mathcal I_\tau(P)=\{\tau(m):m\in\mathcal M,\ P_O(m)=P\}.
\]
Point identification means that this set is a singleton. An empty set signals incompatibility of data and assumptions. Coordinate bounds need not describe the joint set. Bounds are sharp only if every retained target is attainable or arbitrarily approachable by compatible models. Finite bounds require explicit restrictions on unobserved quantities; unrestricted confounding, hidden wealth, extrapolation or arbitrary equilibrium selection can yield uninformative sets.

\textbf{Causal interventions.} A potential outcome \(Y_i(p)\) is the outcome under a complete policy regime \(p\), including timing, financing, information and market spillovers. The target population and units are fixed before treatment. With interference, use the complete assignment vector or a justified exposure mapping; individual no-interference notation alone cannot identify an economy-wide intervention. A difference of mean potential outcomes does not require the cross-world joint distribution. Conditioning on post-treatment migration, survival, enrollment or employment changes the estimand and needs separate justification.

\textbf{Statistical statements.} An estimator, confidence set, forecast or test specifies a sampling law, model class and loss/coverage criterion. Uniform \((1-\alpha)\)-coverage means \(\inf_{m\in\mathcal M}\Pr_m\{\tau(m)\in C_n\}\geq1-\alpha\), or an explicitly asymptotic version. The whole parameter space trivially has coverage, so informativeness requires a stated width, risk or power objective. A forecasting improvement is relative to a named benchmark, information set and evaluation population.

\textbf{Equilibrium and optimization.} An equilibrium comprises feasible plans, optimality, beliefs where needed, budget constraints and all market/resource clearing. The primitives or a stated benchmark define each correspondence \(\mathcal E(p;m)\). Nonemptiness is assumed or proved, never inferred just from compact policy sets. A selected-equilibrium target uses a declared selection rule; a robust target quantifies over all permitted equilibria. Use suprema or infima unless attainment is established. An \(\varepsilon\)-optimal policy is feasible and within \(\varepsilon\) of the finite optimum value. Report an empty feasible set separately.

\textbf{Welfare and accounting.} Normative utility, interpersonal normalization, social weights, numeraire, discounting and the use of public receipts are primitives. Behavioral utility need not coincide with the welfare criterion. Household welfare includes ownership income and rents once; adding profits again to compensating variation generally double-counts them. A monetary equivalent is defined by an explicit transfer and price convention, with monotonicity and range conditions. Average effects, mechanism contributions, transfers and welfare losses are different targets. Infinite sums require convergence and dynamic feasible plans require terminal or no-Ponzi conditions.

\textbf{Source versions.} A working-paper date, revision date and journal publication year are distinguished. A DOI for a journal version is not automatically the DOI of an earlier manuscript. Locators refer to the stated version and printed pages where specified. A metadata-only check does not verify a claim about a full-text conclusion. Every entry's source subsection narrows the provenance claim accordingly.
% END ECONOMICS STATEMENT
\end{document}
